\chapter{State machines and modes}
\label{ch:fsm}

\begin{aside}
A modal app has a small set of modes and explicit transitions
between them. Drawing the modes as a state machine forces you
to think about every transition, not just the happy path.
\end{aside}

\section{The pattern}

\begin{lstlisting}
enum class Mode { Normal, Insert, Visual, Command };
auto mode = signal(Mode::Normal);

void on_key(KeyEvent ev) {
    switch (ctx.get(mode)) {
        case Mode::Normal:  handle_normal(ev);  break;
        case Mode::Insert:  handle_insert(ev);  break;
        // ...
    }
}
\end{lstlisting}

The mode is a signal; each mode has its own dispatcher.

\section{Transitions}

A transition is a mode change plus optional side effects:

\begin{lstlisting}
void enter_insert(int col) {
    move_cursor_to(col);
    ctx.set(mode, Mode::Insert);
}
\end{lstlisting}

Enter and exit hooks fire on transition (e.g., save state
on exit-from-insert).

\section{Hierarchical states}

Some modes have sub-modes. ``Normal mode + waiting for g
prefix'' is a sub-state of normal. Use nested enums or a
state stack.

\section{Visual feedback}

The status bar shows the current mode. Cursor shape
reflects mode (Chapter~\ref{ch:cursor}). The user always
knows where they are.

\section{Mode-specific keymaps}

Each mode has its own keymap. Switching mode swaps the
active keymap.

\section{Diagrams}

For complex state machines, draw the diagram and check
it in. The diagram is the source of truth; the code
implements it.

\section{Statechart libraries}

For very complex apps, statechart libraries (boost.sml,
xstate-cpp) encode the machine declaratively. \maya{}
doesn't bundle one; pick if you need it.

\section{Recap}

\begin{recap}
\begin{itemize}[leftmargin=*]
  \item Mode = signal of an enum.
  \item Each mode has its own dispatcher and keymap.
  \item Enter/exit hooks for side effects.
  \item Hierarchical states for sub-modes.
  \item Visual feedback (status bar, cursor shape).
\end{itemize}
\end{recap}

\section*{Workshop}

\begin{enumerate}[leftmargin=*]
  \item Build a vim-style modal editor.
  \item Add a sub-mode (g-prefix).
  \item Draw the state diagram; check it in.
\end{enumerate}
